\documentclass[11pt,a4paper]{article}

\usepackage{geometry}
\geometry{a4paper, top=2.5cm, bottom=2.5cm, left=2.2cm, right=2.2cm}

\usepackage{fontspec}
\setmainfont{DejaVu Serif}
\setsansfont{DejaVu Sans}
\setmonofont{DejaVu Sans Mono}[Scale=0.90]

\usepackage{amsmath,amssymb}
\usepackage{graphicx}
\graphicspath{{figures/}}

\usepackage{xcolor}
\usepackage{booktabs}
\usepackage{tabularx}
\usepackage{caption}
\usepackage{listings}
\usepackage{fancyhdr}
\usepackage{microtype}
\usepackage{tcolorbox}
\tcbuselibrary{skins,breakable}

\usepackage{hyperref}
\hypersetup{
    colorlinks=true,
    linkcolor=blue!75!black,
    urlcolor=teal!70!black,
    citecolor=blue!75!black,
    pdfauthor={Θάνος},
    pdftitle={Λειτουργικά Συστήματα - Δραστηριότητα 1: Εισαγωγή στον Πυρήνα xv6}
}

% --- Listing Styles ---
\definecolor{codebg}{rgb}{0.97,0.97,0.97}
\definecolor{codeborder}{rgb}{0.85,0.85,0.85}
\definecolor{codegreen}{rgb}{0.1,0.5,0.2}
\definecolor{codegray}{rgb}{0.4,0.4,0.4}
\definecolor{codepurple}{rgb}{0.5,0.1,0.7}
\definecolor{codeblue}{rgb}{0.1,0.2,0.8}

\lstdefinestyle{cstyle}{
    backgroundcolor=\color{codebg},
    commentstyle=\color{codegreen}\itshape,
    keywordstyle=\color{codeblue}\bfseries,
    numberstyle=\tiny\color{codegray},
    stringstyle=\color{codepurple},
    basicstyle=\ttfamily\footnotesize,
    breakatwhitespace=false,
    breaklines=true,
    captionpos=b,
    keepspaces=true,
    numbers=left,
    numbersep=7pt,
    showspaces=false,
    showstringspaces=false,
    showtabs=false,
    tabsize=4,
    frame=single,
    rulecolor=\color{codeborder}
}

\lstdefinestyle{termstyle}{
    backgroundcolor=\color{black!92},
    basicstyle=\ttfamily\footnotesize\color{white},
    breaklines=true,
    frame=single,
    rulecolor=\color{black},
    numbers=none
}

\lstset{style=cstyle}

% --- Problem Prompt Box ---
\newtcolorbox{promptbox}[1][]{
    colback=blue!3!white,
    colframe=blue!40!black,
    fonttitle=\bfseries,
    coltitle=white,
    title=Εκφώνηση,
    arc=1mm,
    boxrule=0.6pt,
    left=3mm, right=3mm, top=2mm, bottom=2mm,
    breakable,
    #1
}

% --- Header / Footer ---
\pagestyle{fancy}
\fancyhf{}
\fancyhead[L]{\small\textsc{Λειτουργικά Συστήματα (ECE ΓΚ702)}}
\fancyhead[R]{\small\textsc{Δραστηριότητα 1}}
\fancyfoot[C]{\thepage}
\renewcommand{\headrulewidth}{0.4pt}

\begin{document}

% --- Header Block ---
\begin{center}
    {\large \textbf{ΠΑΝΕΠΙΣΤΗΜΙΟ ΠΑΤΡΩΝ}}\\[1mm]
    {\normalsize Τμήμα Ηλεκτρολόγων Μηχανικών \& Τεχνολογίας Υπολογιστών}\\[3mm]
    {\large \textbf{Λειτουργικά Συστήματα (ECE ΓΚ702)}}\\[1mm]
    \rule{\textwidth}{0.8pt}\\[3mm]
    {\LARGE \textbf{Δραστηριότητα 1: Εισαγωγή στον Πυρήνα xv6 \& System Calls}}\\[2mm]
    \rule{\textwidth}{0.8pt}\\[3mm]
    \begin{tabular}{ll}
        \textbf{Φοιτητής:} & [Παναγόπουλος Αθανάσιος] \\
        \textbf{Αριθμός Μητρώου (ΑΜ):} & [1107645] \\
        \textbf{Ημερομηνία:} & \today \\
        \textbf{Repository:} & \url{https://github.com/RealThanosP/thanOS-v6} \\
    \end{tabular}
\end{center}
\vspace{2mm}


% ==============================================================================
% ΑΣΚΗΣΗ 1
% ==============================================================================
\section{Άσκηση 1: Εντοπισμός System Calls \& Ανάλυση Προγράμματος Χρήστη}


\begin{promptbox}
Βρείτε το αρχείο στον πηγαίο κώδικα του xv6 όπου ορίζονται οι κλήσεις συστήματος οι οποίες θα
είναι διαθέσιμες σε προγράμματα χρήστη. Έπειτα, διαλέξτε ένα από τα προγράμματα χρήστη στον
φάκελο \texttt{user}. Εξηγήστε τη λειτουργία του κάνοντας συγκεκριμένες αναφορές στα system calls που
καλεί.
\end{promptbox}

\subsection*{1Α. Αρχεία Ορισμού Κλήσεων Συστήματος}

Στον πηγαίο κώδικα του xv6, ο ορισμός και η διεπαφή (interface) των κλήσεων συστήματος (system calls) για τα προγράμματα χρήστη κατανέμονται σε τρία βασικά αρχεία:

\begin{itemize}\setlength{\itemsep}{1pt}\setlength{\parskip}{1pt}
    \item \textbf{\texttt{include/user.h}}: Περιέχει τις δηλώσεις των συναρτήσεων (function prototypes) των διαθέσιμων κλήσεων συστήματος για το user space (π.χ. \texttt{int open(const char*, int)}, \texttt{int read(int, void*, int)}, \texttt{int write(int, const void*, int)}, \texttt{int close(int)}, \texttt{int exit(void)}, κ.λπ.). Με τον τρόπο αυτό, ο μεταγλωττιστής γνωρίζει τους τύπους ορισμάτων και επιστροφής κατά τη συγγραφή προγραμμάτων χρήστη.
    
    \item \textbf{\texttt{include/syscall.h}}: Ορίζει τους αριθμητικούς κωδικούς (system call numbers) μέσω macros του προεπεξεργαστή (π.χ. \texttt{\#define SYS\_open 15}, \texttt{\#define SYS\_read 5}, \texttt{\#define SYS\_write 16}, \texttt{\#define SYS\_close 21}, \texttt{\#define SYS\_exit 2}). Οι κωδικοί αυτοί χρησιμοποιούνται τόσο από το user space για την επιλογή της κλήσης, όσο και από τον πυρήνα για την ανάκτηση της κατάλληλης συνάρτησης χειρισμού από τον πίνακα δεικτών \texttt{syscalls[]}.
    
    \item \textbf{\texttt{ulib/usys.S}}: Περιέχει τα assembly wrappers (stubs) που εκτελούνται στο user space. Μέσω του macro \texttt{SYSCALL(name)}:
\begin{lstlisting}[language=C, aboveskip=2pt, belowskip=2pt, basicstyle=\ttfamily\scriptsize]
#define SYSCALL(name) \
  .globl name; name: movl $SYS_ ## name, %eax; int $T_SYSCALL; ret
\end{lstlisting}
    φορτώνεται ο αντίστοιχος κωδικός στον καταχωρητή \texttt{\%eax} και εκτελείται η διακοπή λογισμικού \texttt{int \$T\_SYSCALL} (\texttt{trap 64}), μεταφέροντας τον έλεγχο από το User Space (Ring 3) στον Πυρήνα (Ring 0).
\end{itemize}

\clearpage
\subsection*{1Β. Επιλογή \& Ανάλυση Προγράμματος Χρήστη (\texttt{user/cat.c})}

Για την ανάλυση επιλέγουμε το πρόγραμμα χρήστη \texttt{cat} (\texttt{user/cat.c}), το οποίο αναλαμβάνει το άνοιγμα ενός ή περισσότερων αρχείων και την εμφάνιση των περιεχομένων τους στο τερματικό.

% %%%%%%%%%%%%%%%%%%%%%%%%%%%%%%5
% Source code of cat.c
\subsubsection*{Πηγαίος Κώδικας (\texttt{user/cat.c})}
\begin{lstlisting}[language=C, caption={Ο πηγαίος κώδικας του προγράμματος cat}]
#include "types.h"
#include "stat.h"
#include "user.h"

char buf[512];

void
cat(int fd)
{
  int n;

  while((n = read(fd, buf, sizeof(buf))) > 0)
    write(1, buf, n);
  if(n < 0){
    printf(1, "cat: read error\n");
    exit();
  }
}

int
main(int argc, char *argv[])
{
  int fd, i;

  if(argc <= 1){
    cat(0);
    exit();
  }

  for(i = 1; i < argc; i++){
    if((fd = open(argv[i], 0)) < 0){
      printf(1, "cat: cannot open %s\n", argv[i]);
      exit();
    }
    cat(fd);
    close(fd);
  }
  exit();
}
\end{lstlisting}
% ===============

\clearpage

\subsubsection*{Ανάλυση Ροής Εκτέλεσης \& System Calls}

\begin{enumerate} 
    \setlength{\itemsep}{2pt}
    \setlength{\parskip}{1pt}

    \item \textbf{Διαχείριση Ορισμάτων \& Είσοδος}:
    \begin{itemize}\setlength{\itemsep}{1pt}\setlength{\parskip}{0pt}
        \item Αν δεν έχουν δοθεί ορίσματα αρχείων στη γραμμή εντολών (\texttt{argc <= 1}), το πρόγραμμα καλεί \texttt{cat(0)}, διαβάζοντας απευθείας από την πρότυπη είσοδο (standard input, file descriptor 0), και έπειτα τερματίζει με \texttt{exit()}.
        \item Αν δοθούν ονόματα αρχείων (\texttt{argc > 1}), το \texttt{cat} επεξεργάζεται σειριακά κάθε αρχείο \texttt{argv[i]} μέσω ενός βρόχου \texttt{for}.
    \end{itemize}

    \item \textbf{Άνοιγμα Αρχείου (\texttt{SYS\_open})}:
    \begin{itemize}
        \item Για κάθε αρχείο εκτελείται η κλήση συστήματος \texttt{open(argv[i], 0)} (\texttt{SYS\_open}, κωδικός 15) με flag ανάγνωσης (\texttt{O\_RDONLY}).
        \item Εάν η επιστρεφόμενη τιμή είναι αρνητική (\texttt{fd < 0}), σημαίνει ότι το αρχείο δεν μπόρεσε να ανοιχθεί. Το πρόγραμμα τυπώνει μήνυμα σφάλματος και τερματίζει καλώντας \texttt{exit()} (\texttt{SYS\_exit}, κωδικός 2).
    \end{itemize}

    \item \textbf{Ανάγνωση Περιεχομένου (\texttt{SYS\_read})}:
    \begin{itemize}
        \item Μέσα στη συνάρτηση \texttt{cat(fd)}, εκτελείται επαναληπτικά η κλήση συστήματος \texttt{read(fd, buf, sizeof(buf))} (\texttt{SYS\_read}, κωδικός 5) για να διαβάσει τμήματα δεδομένων μεγέθους έως 512 bytes στον προσωρινό buffer \texttt{buf}.
        \item Η τιμή επιστροφής \texttt{n} αντιστοιχεί στα πραγματικά bytes που διαβάστηκαν.
    \end{itemize}

    \item \textbf{Εγγραφή στην Έξοδο (\texttt{SYS\_write})}:
    \begin{itemize}
        \item Για όσο \texttt{n > 0}, καλείται η \texttt{write(1, buf, n)} (\texttt{SYS\_write}, κωδικός 16), στέλνοντας τα δεδομένα του buffer στον file descriptor 1 (standard output), που αντιστοιχεί στην οθόνη του τερματικού.
    \end{itemize}

    \item \textbf{Τερματισμός Βρόχου Ανάγνωσης}:
    \begin{itemize}
        \item Τα βήματα 3 και 4 εκτελούνται διαδοχικά έως ότου η \texttt{read()} επιστρέψει \texttt{0}, γεγονός που σημαίνει ότι επιτεύχθηκε το τέλος του αρχείου (End-Of-File - EOF).
        \item Αν η \texttt{read()} επιστρέψει αρνητική τιμή (\texttt{n < 0}), έχει συμβεί σφάλμα ανάγνωσης, τυπώνεται σφάλμα και καλείται \texttt{exit()} (\texttt{SYS\_exit}, κωδικός 2).
    \end{itemize}

    \item \textbf{Κλείσιμο Περιγραφέα (\texttt{SYS\_close})}:
    \begin{itemize}
        \item Μόλις διαβαστεί πλήρως το αρχείο, καλείται η \texttt{close(fd)} (\texttt{SYS\_close}, κωδικός 21), απελευθερώνοντας τον περιγραφέα αρχείου και τους δεσμευμένους πόρους του kernel.
    \end{itemize}

    \item \textbf{Τερματισμός Διεργασίας (\texttt{SYS\_exit})}:
    \begin{itemize}
        \item Αφού ολοκληρωθεί η επεξεργασία όλων των αρχείων, καλείται η \texttt{exit()} (\texttt{SYS\_exit}, κωδικός 2), η οποία τερματίζει τη διεργασία \texttt{cat} και επιστρέφει τον έλεγχο στο shell.
    \end{itemize}
\end{enumerate}

\clearpage


% ==============================================================================
% ΑΣΚΗΣΗ 2
% ==============================================================================
\section{Άσκηση 2: Δημιουργία Προγράμματος Χρήστη}

\begin{promptbox}
Δημιουργήστε ένα πρόγραμμα χρήστη για το xv6.
Μέσα στο φάκελο \texttt{user}, δημιουργήστε ένα αρχείο με ό,τι όνομα θέλετε. Εκεί, γράψτε ένα πρόγραμμα
σε γλώσσα C το οποίο θα κάνει κάτι χρήσιμο. Δώστε προσοχή σε ποιες συναρτήσεις έχετε πρόσβαση. Εδώ,
δεν μπορείτε να χρησιμοποιήσετε τη βασική βιβλιοθήκη της C που ξέρετε αλλά τη βασική βιβλιοθήκη
που έχει δημιουργηθεί για το xv6 (\texttt{ulib.c}).
Έπειτα, επεξεργαστείτε κατάλληλα το αρχείο \texttt{Makefile} ώστε να προστεθεί το πρόγραμμα στο image
του xv6 κατά την κατασκευή του (προσθήκη στη λίστα \texttt{USER\_PROGS}).
Εκτελέστε \texttt{make qemu} και χρησιμοποιήστε το πρόγραμμα που δημιουργήσατε μέσα στο xv6.
\end{promptbox}

\subsection*{2A. Περιγραφή \& Σκοπός Προγράμματος}
% TODO: Περιγράψτε το πρόγραμμα που υλοποιήσατε (π.χ. touch) και τον τρόπο χρήσης του.
Υλοποίησα την εντολή \texttt{touch} η οποία δημιουργεί ένα νέο αρχείο με (Read/Write permissions) με το όνομα που της δίνεται ως όρισμα. Αν το αρχείο υπάρχει ήδη, τότε δέν γίνεται τίποτα. 
Η χρήση της εντολής είναι απλή: \texttt{touch <filename>}.

\subsection*{2B. Πηγαίος Κώδικας (\texttt{user/touch.c})}
\begin{lstlisting}[language=C, caption={Υλοποίηση προγράμματος χρήστη}]
#include "types.h"
#include "stat.h"
#include "user.h"
#include "fcntl.h"

int touch()
{
    // Return flags for creating a new file with READ/WRITE permissions
    // Creates a file if not exists with O_CREATE and gives it read/write access with O_RDWR flag
    return O_CREATE | O_RDWR;
}

int main(int argc, char *argv[])
{
    int fd, i;

    if (argc < 2)
    {
        printf(2, "Usage: touch file1 file2 ...\n");
        exit();
    }

    for (i = 1; i < argc; i++)
    {
        fd = open(argv[i], touch());
        if (fd < 0)
        {
            printf(2, "touch: cannot create %s\n", argv[i]);
            exit();
        }
        close(fd);
    }
    exit();
}
\end{lstlisting}

\subsection*{2C. Ενσωμάτωση στο \texttt{Makefile}}
\begin{lstlisting}[language=make, caption={Προσθήκη στη μεταβλητή USER\_PROGS}]
USER_PROGS=\
# ...
	touch\
# ...
\end{lstlisting}
Εφόσον δεν προσθέσαμε κάποιο νέο αρχείο βιβλιοθήκης ή \texttt{system call}, δεν χρειάστηκε να
τροποποιήσουμε κάτι άλλο. 

\subsection*{2D. Έλεγχος \& Έξοδος Εκτέλεσης στο QEMU}
\begin{lstlisting}[style=termstyle, caption={Εκτέλεση στο τερματικό του xv6}]
$ touch hello
$ ls
.              1 1 512
..             1 1 512
README         2 2 1588
cat            2 3 15512
echo           2 4 14752
forktest       2 5 9712
schedtest      2 6 17696
grep           2 7 17400
init           2 8 15304
kill           2 9 14848
ln             2 10 14808
ls             2 11 17176
mkdir          2 12 14896
rm             2 13 14880
sh             2 14 25792
stressfs       2 15 15504
usertests      2 16 50752
wc             2 17 15920
zombie         2 18 14536
lotteryschedte 2 19 18120
touch          2 20 15208
shutdown       2 21 14664
favnum         2 22 14656
syscount       2 23 15016
killrandom     2 24 14672
clear          2 25 14608
console        3 26 0
hello          2 27 0 // GOOD
\end{lstlisting}


% ==============================================================================
% ΑΣΚΗΣΗ 3
% ==============================================================================
\section{Άσκηση 3: Ανάλυση της Κλήσης Συστήματος \texttt{uptime}}

\begin{promptbox}
Αναζητώντας μέσα στον πηγαίο κώδικα του xv6, βρείτε την κλήση συστήματος \texttt{uptime}.
\begin{enumerate}
    \item Βρείτε όλα τα σημεία του κώδικα του πυρήνα στα οποία γίνεται αναφορά σε αυτή.
    \item Εξηγήστε την υλοποίησή της.
\end{enumerate}
\end{promptbox}

\subsection*{1. Σημεία Αναφοράς στον Πηγαίο Κώδικα}
% TODO: Καταγράψτε όλα τα αρχεία και γραμμές (π.χ. syscall.h, syscall.c, sysproc.c, user.h, usys.S).
Το \texttt{uptime} ορίζεται ως κλήση συστήματος στο kernel του xv6 και η αναφορά της εμφανίζεται στα εξής αρχεία:
\begin{itemize}
    \item \texttt{include/syscall.h}: Δηλώνεται ο αριθμός της.
    \item \texttt{include/user.h}: Δηλώνεται η συνάρτηση που είναι διαθέσιμη για χρήση στο \texttt{user space}.
    \item \texttt{kernel/syscall.c}: Δηλώνεται η συνάρτησή υλοποίσης της και προστίθεται στην λίστα των διαθέσιμων \texttt{system call} του \texttt{xv6}.
    \item \texttt{kernel/sysproc.h}: Ορίζεται η συνάρτηση \texttt{sys\_uptime()} που την υλοποιεί.
    \item \texttt{ulib/usys.S}: Δηλώνεται ώς διαθέσιμο system call στον assembly wrapper.
\end{itemize}

\subsection*{2. Ανάλυση Υλοποίησης (\texttt{sys\_uptime})}
% TODO: Παραθέστε τον κώδικα της συνάρτησης και εξηγήστε:
% - Τι μετράει η μεταβλητή ticks (hardware timer interrupts).
% - Γιατί χρησιμοποιείται το lock (tickslock).
Επιστρέφει τον αριθμό των διακοπών ρολογιού (clock tick interrupts) που έχουν συμβεί από την εκκίνηση του συστήματος. 
Η μεταβλητή \texttt{ticks} αυξάνεται κάθε φορά που συμβαίνει μια διακοπή ρολογιού (\texttt{kernel/trap.c}).
Ο λόγος που χρησιμοποιείται το \texttt{tickslock} είναι για να αποφευχθεί η ταυτόχρονη πρόσβαση στη μεταβλητή \texttt{ticks} από πολλαπλά threads ή διεργασίες, εξασφαλίζοντας ότι η ανάγνωση της τιμής της είναι ασφαλής και συνεπής.
Στην ουσία παγώνουμε τον CPU και εξασφαλίζουμε ότι μόνο η uptime θα έχει πρόσβαση στην \texttt{ticks} για την εκτέλεση.

\begin{lstlisting}[language=C, caption={Υλοποίηση sys\_uptime στο kernel/sysproc.c}]
// return how many clock tick interrupts have occurred
// since start.
int sys_uptime(void)
{
  uint xticks;

  acquire(&tickslock);
  xticks = ticks;
  release(&tickslock);
  return xticks;
}
\end{lstlisting}

% ==============================================================================
% ΑΣΚΗΣΗ 4
% ==============================================================================
\section{Άσκηση 4: Ροή Εκτέλεσης System Calls στον Πυρήνα}

\begin{promptbox}
Βρείτε τη συνάρτηση μέσα στον kernel κώδικα όπου τελικά εκτελούνται οι system calls. Καταγράψτε
τη σειρά των ενεργειών που γίνονται όταν καλείται ένα system call.
\end{promptbox}

\subsection*{1. Εντοπισμός Συνάρτησης Εκτέλεσης System Calls}
% TODO: Αναφορά στο αρχείο kernel/syscall.c και στη συνάρτηση syscall(void).
H συνάρτηση που εκτελεί όλα τα \texttt{system calls} είναι

\subsection*{2. Σειρά Ενεργειών \& Μηχανισμός Εκτέλεσης System Calls}
% TODO: Εξηγήστε τα βήματα:
% 1. Πώς μεταβιβάζεται ο αριθμός της κλήσης (μέσω eax στο trap frame).
% 2. Έλεγχος ορίων (bounds checking).
% 3. Κλήση της αντίστοιχης συνάρτησης από τον πίνακα δεικτών συναρτήσεων syscalls[].
% 4. Αποθήκευση της τιμής επιστροφής στο eax.


% ==============================================================================
% ΑΣΚΗΣΗ 5
% ==============================================================================
\section{Άσκηση 5: Προσθήκη System Call (\texttt{getfavnum})}

\begin{promptbox}
Προσθέστε μία κλήση συστήματος η οποία θα επιστρέφει τον αγαπημένο σας αριθμό. Η κλήση
που θα φτιάξετε πρέπει να έχει το εξής πρωτότυπο: \texttt{int getfavnum(void)}. Μία τέτοια κλήση δεν
έχει κανέναν πρακτικό σκοπό αφού δεν εκμεταλλεύεται κάποια πληροφορία ή πόρο που κατέχει το
λειτουργικό σύστημα. Στόχος εδώ είναι να κατανοήσουμε τον τρόπο με τον οποίο προστίθενται
κλήσεις συστήματος στο xv6.
\end{promptbox}

\subsection*{1. Τροποποιήσεις στον Πυρήνα \& User Interface}
% TODO: Καταγράψτε τα αρχεία που τροποποιήθηκαν:
% - include/syscall.h
% - ulib/usys.S
% - include/user.h
% - kernel/syscall.c
% - kernel/sysproc.c

\subsection*{2. Πρόγραμμα Ελέγχου (\texttt{user/favnum.c})}
\begin{lstlisting}[language=C, caption={Πρόγραμμα δοκιμής favnum}]
// TODO: Κώδικας προγράμματος χρήστη
\end{lstlisting}

\subsection*{3. Αποτελέσματα Εκτέλεσης}
\begin{lstlisting}[style=termstyle]
$ # TODO: Logs εκτέλεσης
\end{lstlisting}


% ==============================================================================
% ΑΣΚΗΣΗ 6
% ==============================================================================
\section{Άσκηση 6: System Call Τερματισμού (\texttt{halt} / \texttt{shutdown})}

\begin{promptbox}
Προσθέστε μία κλήση συστήματος η οποία θα εκτελεί λειτουργία «shutdown». Η κλήση που θα
φτιάξετε πρέπει να έχει το εξής πρωτότυπο: \texttt{void halt(void)}. Έπειτα, φτιάξτε ένα πρόγραμμα
χρήστη με όνομα \texttt{shutdown} το οποίο θα καλεί την κλήση συστήματος που φτιάξατε, ώστε χρήστες
του xv6 να μπορούν πλέον να κάνουν shutdown το σύστημα χωρίς να χρειάζεται να το κάνουν από
το host μηχάνημα.

Σε αυτή τη κλήση συστήματος θα πρέπει να επικοινωνήσετε άμεσα με το hardware της εικονικής
μηχανής (στην περίπτωσή μας το QEMU). Για την απενεργοποίηση του συστήματος το hardware περιμένει
σε κάποια συγκεκριμένη διεύθυνση (port) να υπάρξει μία συγκεκριμένη τιμή. Εξηγήστε τη λειτουργία
της συνάρτησης \texttt{outw} την οποία θα χρησιμοποιήσετε.
\end{promptbox}

\subsection*{1. Ανάλυση της Συνάρτησης \texttt{outw} \& Port I/O}
% TODO: Εξηγήστε:
% - Port-mapped I/O (x86 in/out εντολές) vs Memory-mapped I/O.
% - Προνομιακές εντολές (Ring 0 vs Ring 3 / GP Fault).
% - ACPI I/O Ports και magic values (0x604, 0x2000 / PM1a_CNT).

\subsection*{2. Υλοποίηση Κλήσης Συστήματος (\texttt{kernel/sysproc.c})}
\begin{lstlisting}[language=C, caption={Υλοποίηση sys\_halt}]
// TODO: Κώδικας sys_halt
\end{lstlisting}

\subsection*{3. Πρόγραμμα Χρήστη \texttt{shutdown} \& Επαλήθευση}
\begin{lstlisting}[language=C, caption={user/shutdown.c}]
// TODO: Κώδικας shutdown.c
\end{lstlisting}


% ==============================================================================
% ΑΣΚΗΣΗ 7
% ==============================================================================
\section{Άσκηση 7: Καταμέτρηση Κλήσεων Συστήματος (\texttt{getcount})}

\begin{promptbox}
Προσθέστε μία κλήση συστήματος η οποία θα επιστρέφει το πόσες φορές έχει εκτελεστεί μία
συγκεκριμένη κλήση συστήματος η οποία θα περνιέται σαν όρισμα. Η κλήση που θα φτιάξετε πρέπει
να έχει το εξής πρωτότυπο: \texttt{int getcount(int syscall)}.
\end{promptbox}

\subsection*{1. Αρχιτεκτονική Καταμέτρησης στον Πυρήνα}
% TODO: Περιγράψτε:
% - Πίνακας καταμέτρησης στο kernel/syscall.c
% - Πού και πώς αυξάνεται ο μετρητής (μέσα στη συνάρτηση syscall())
% - Έλεγχος εγκυρότητας του ορίσματος syscall

\subsection*{2. Υλοποίηση \texttt{sys\_getcount}}
\begin{lstlisting}[language=C, caption={Υλοποίηση sys\_getcount στο kernel/syscall.c}]
// TODO: Κώδικας sys_getcount
\end{lstlisting}

\subsection*{3. Πρόγραμμα Ελέγχου \& Πειραματική Επαλήθευση}
\begin{lstlisting}[style=termstyle, caption={Δοκιμές καταμέτρησης}]
$ # TODO: Terminal logs δοκιμών
\end{lstlisting}


% ==============================================================================
% ΑΣΚΗΣΗ 8
% ==============================================================================
\section{Άσκηση 8: Τυχαίος Τερματισμός Διεργασίας (\texttt{killrandom})}

\begin{promptbox}
Προσθέστε μία κλήση συστήματος η οποία θα τερματίζει μια τυχαία διεργασία με πρότυπο
\texttt{int killrandom(void)}.
Θα χρειαστεί να προσθέσετε μία γεννήτρια τυχαίων αριθμών στο kernel. Πώς μπορείτε να βρείτε
τα PID όλων των τρέχοντων διεργασιών στο σύστημα;
\end{promptbox}

\subsection*{1. Γεννήτρια Ψευδοτυχαίων Αριθμών (PRNG) στον Πυρήνα}
% TODO: Περιγράψτε τον αλγόριθμο PRNG (π.χ. Xorshift32 με βάση τα ticks) και την υλοποίησή του.

\subsection*{2. Ανάκτηση PID \& Προστασία Κρίσιμων Διεργασιών}
% TODO: Εξηγήστε:
% - Πώς διατρέχεται ο πίνακας διεργασιών (ptable.proc)
% - Πώς προστατεύεται η διεργασία init (PID 1) και η τρέχουσα διεργασία (proc)
% - Συγχρονισμός με ptable.lock

\subsection*{3. Υλοποίηση \texttt{sys\_killrandom}}
\begin{lstlisting}[language=C, caption={Υλοποίηση sys\_killrandom στο kernel/sysproc.c}]
// TODO: Κώδικας sys_killrandom
\end{lstlisting}

\subsection*{4. Πρόγραμμα Χρήστη \& Δοκιμή}
\begin{lstlisting}[style=termstyle]
$ # TODO: Logs τερματισμού τυχαίων διεργασιών
\end{lstlisting}

\end{document}
